\section{August 26, 2026}

The main topics are properties of solutions, time discretization, and
consistency error. We continue to consider the IVP
\begin{equation}\label{eq:ivp-aug-26}
  x'(t)=f(t,x(t)),
  \qquad
  x(t_0)=x_0,
\end{equation}
where $x\colon \bb{R}\to\bb{R}^n$.

\subsection{Stability of solutions}

The integral formulation of \eqref{eq:ivp-aug-26} is
\[
  x(t)=x_0+\int_{t_0}^t f(s,x(s))\,ds.
\]
If $y(t)$ is another solution with initial value $y_0$, then
\[
  y(t)=y_0+\int_{t_0}^t f(s,y(s))\,ds,
\]
and hence
\[
  y(t)-x(t)
  =y_0-x_0
   +\int_{t_0}^t\bigl(f(s,y(s))-f(s,x(s))\bigr)\,ds.
\]
Suppose $f$ is Lipschitz continuous in $x$ with Lipschitz constant $L$,
so that
\[
  \lVert f(t,y)-f(t,x)\rVert\leq L\lVert y-x\rVert.
\]
Taking norms gives
\begin{equation}\label{eq:solution-difference-aug-26}
  \lVert y(t)-x(t)\rVert
  \leq \lVert y_0-x_0\rVert
  +\int_{t_0}^t L\lVert y(s)-x(s)\rVert\,ds.
\end{equation}
The following inequality controls this accumulation of error.

\begin{lemma}[Gr\"onwall's inequality]\label{lem:gronwall}
  Let $u,v\in C([t_0,t_1],\bb{R})$ be nonnegative functions and let
  $\rho\geq 0$. If
  \begin{equation}\label{eq:gronwall-hypothesis}
    u(t)\leq \rho+\int_{t_0}^t v(s)u(s)\,ds,
    \qquad t\in[t_0,t_1],
  \end{equation}
  then
  \begin{equation}\label{eq:gronwall-conclusion}
    u(t)\leq \rho
    \exp\left(\int_{t_0}^t v(s)\,ds\right),
    \qquad t\in[t_0,t_1].
  \end{equation}
  In particular, $u\equiv 0$ when $\rho=0$.
\end{lemma}

\begin{proof}
  First suppose that $\rho>0$ and define
  \[
    U(t)=\rho+\int_{t_0}^t v(s)u(s)\,ds.
  \]
  Then $U\in C^1([t_0,t_1],\bb{R})$, and
  \[
    u(t)\leq U(t),
    \qquad
    U(t)\geq \rho>0.
  \]
  Since $u$ and $v$ are nonnegative,
  \[
    U'(t)=v(t)u(t)\leq v(t)U(t).
  \]
  Thus the logarithmic derivative of $U$ satisfies
  \[
    \bigl(\log U(t)\bigr)'=\frac{U'(t)}{U(t)}\leq v(t).
  \]
  Integrating from $t_0$ to $t$ yields
  \[
    \log U(t)-\log U(t_0)
    \leq \int_{t_0}^t v(s)\,ds.
  \]
  Because $U(t_0)=\rho$ and $u(t)\leq U(t)$, applying the exponential
  function proves \eqref{eq:gronwall-conclusion} when $\rho>0$.

  Now suppose that $\rho=0$. For any $\varepsilon>0$,
  \eqref{eq:gronwall-hypothesis} implies
  \[
    u(t)\leq \varepsilon+\int_{t_0}^t v(s)u(s)\,ds.
  \]
  The first part of the proof, with $\rho=\varepsilon$, gives
  \[
    0\leq u(t)
    \leq \varepsilon
    \exp\left(\int_{t_0}^t v(s)\,ds\right).
  \]
  Letting $\varepsilon\downarrow 0$ shows that $u(t)=0$ for every
  $t\in[t_0,t_1]$.
\end{proof}

Apply Lemma~\ref{lem:gronwall} to
\eqref{eq:solution-difference-aug-26} with
\[
  u(t)=\lVert y(t)-x(t)\rVert,
  \qquad
  v(t)=L,
  \qquad
  \rho=\lVert y_0-x_0\rVert.
\]
We obtain
\begin{equation}\label{eq:gronwall-stability-aug-26}
  \lVert y(t)-x(t)\rVert
  \leq \lVert y_0-x_0\rVert e^{L(t-t_0)}.
\end{equation}
This estimate can be pessimistic, but it establishes continuous
dependence on the initial data.

\subsection{The solution map}

The \emph{solution map}, or \emph{flow map}, is denoted by
\[
  x(t)=\varphi^{t,t_0}(x_0).
\]
For example, consider the scalar logistic equation
\[
  x'=x(1-x).
\]
For $x_0\neq 0$, its solution is
\[
  \varphi^{t,t_0}(x_0)
  =\frac{1}{1+\dfrac{1-x_0}{x_0}e^{-(t-t_0)}}.
\]
When $x_0=0$, the solution is $x(t)\equiv0$.

The solution map has the identity and composition properties
\begin{align*}
  \varphi^{t_0,t_0}(x_0)&=x_0,
  &\varphi^{t_0,t_0}&=I,\\
  \varphi^{t_2,t_1}\bigl(\varphi^{t_1,t_0}(x_0)\bigr)
    &=\varphi^{t_2,t_0}(x_0).
\end{align*}
The second property is the basis of time marching.

Differentiating the solution map with respect to time gives
\begin{align*}
  \frac{d}{dt}\varphi^{t,t_0}(x_0)
    &=f(t,x(t)),\\
  \frac{d^2}{dt^2}\varphi^{t,t_0}(x_0)
    &=f_t(t,x(t))+f_x(t,x(t))x'(t)\\
    &=\bigl(f_t+f_xf\bigr)(t,x(t)),
\end{align*}
where $f_x$ is the Jacobian matrix of $f$ with respect to $x$. Higher
time derivatives can be found by repeated differentiation, for example
with automatic differentiation; they can also be viewed through the
Koopman operator.

To measure sensitivity with respect to the initial value, define the
Wronskian matrix
\[
  W(t)=\frac{\partial}{\partial x_0}\varphi^{t,t_0}(x_0)
      =\frac{\partial x(t)}{\partial x_0}.
\]
Differentiating with respect to $t$ gives
\begin{align*}
  W'(t)
  &=\frac{\partial}{\partial x_0}f(t,x(t))\\
  &=f_x(t,x(t))\frac{\partial x(t)}{\partial x_0}.
\end{align*}
Thus, if
\[
  A(t)=f_x(t,x(t)),
\]
then $W$ satisfies the \emph{variational equation}
\begin{equation}\label{eq:variational-equation}
  W'(t)=A(t)W(t),
  \qquad
  W(t_0)=I.
\end{equation}

\subsection{One-step methods}

Choose a time grid
\[
  t_0<t_1<\cdots<t_N=t,
  \qquad
  \tau_j=t_{j+1}-t_j,
\]
where the step sizes $\tau_j$ are small. Repeated use of the composition
property gives
\[
  \varphi^{t_N,t_0}
  =\varphi^{t_N,t_{N-1}}\circ\cdots\circ
   \varphi^{t_2,t_1}\circ\varphi^{t_1,t_0}.
\]
It is therefore enough to approximate a single short-time solution map.
Writing $t_j=t$ and $t_{j+1}=t+\tau$, Taylor expansion gives
\begin{equation}\label{eq:flow-taylor-expansion}
  \varphi^{t+\tau,t}(x)
  =x+\tau f(t,x)
   +\frac{\tau^2}{2}\bigl(f_t+f_xf\bigr)(t,x)
   +O(\tau^3).
\end{equation}

Dropping the terms of order $\tau^2$ and higher defines the Euler
one-step map
\[
  \Psi^{t+\tau,t}(x)=x+\tau f(t,x).
\]
Consequently, Euler's method is
\begin{equation}\label{eq:euler-method}
  x_{j+1}=x_j+\tau_j f(t_j,x_j).
\end{equation}
Its one-step consistency error follows from
\eqref{eq:flow-taylor-expansion}:
\[
  \varphi^{t+\tau,t}(x)-\Psi^{t+\tau,t}(x)
  =\frac{\tau^2}{2}\bigl(f_t+f_xf\bigr)(t,x)+O(\tau^3)
  =O(\tau^2).
\]
